\documentclass[11pt,a4paper]{article}

% --- Packages ---
\usepackage[utf8]{inputenc}
\usepackage[T1]{fontenc}
\usepackage{amsmath, amsfonts, amssymb, amsthm}
\usepackage{bm} % Bold math
\usepackage{graphicx}
\usepackage{hyperref}
\usepackage[margin=1in]{geometry}
\usepackage{authblk} % For authors and affiliations
\usepackage{microtype}
\usepackage{booktabs}
\usepackage{cleveref} % Smart referencing
\usepackage[title,titletoc]{appendix}
\usepackage[backend=biber, style=numeric]{biblatex}
\addbibresource{references.bib}

% --- Theorem Environments ---
\newtheorem{theorem}{Theorem}[section]
\newtheorem{lemma}[theorem]{Lemma}
\newtheorem{proposition}[theorem]{Proposition}
\newtheorem{corollary}[theorem]{Corollary}

\theoremstyle{definition}
\newtheorem{definition}[theorem]{Definition}
\newtheorem{example}[theorem]{Example}
\newtheorem{remark}[theorem]{Remark}
\newtheorem{assumption}[theorem]{Assumption}

\numberwithin{equation}{section}

% --- Custom Commands ---
\newcommand{\R}{\mathbb{R}}
\newcommand{\N}{\mathbb{N}}
\newcommand{\E}{\mathbb{E}}
\newcommand{\prob}{\mathbb{P}}
\newcommand{\diff}{\mathrm{d}}

% --- Metadata ---
\title{PASSES-HGV: Expanding the Continuous Spectral Framework to 6-DOF Hypersonic Glide and Fractional Orbital Trajectories}

\author[1]{Ethan Knox\thanks{Email: \texttt{ethank5149@gmail.com}}}
\affil[1]{Independent Researcher, Champaign, IL, USA}

\date{\today}

\begin{document}

\maketitle

\begin{abstract}
This paper presents PASSES-HGV, an open-source, continuous mathematical framework that expands the original fixed-grid spectral formulation into a full 6-Degree-of-Freedom (6-DOF) hypersonic glide and Fractional Orbital Bombardment System (FOBS) architecture. By replacing traditional discrete finite element methods and decoupled aerodynamic look-up tables with a 2D bivariate Chebyshev tensor-product space and Kronecker product operators, the framework maintains $C^\infty$ continuity across complex waverider geometries. Anisotropic flexural rigidity and localized torsional flutter are governed continuously across the planform, while a coupled, unclassified aerothermodynamic model links modified Newtonian pressure and Sutton-Graves convective heating directly to a smoothed enthalpy-based ablation matrix. To ensure autonomous trajectory optimization, an embedded numerical Predictor-Corrector guidance algorithm utilizes a novel exterior penalty function mapped directly into the downrange error space, forcing the shadow solver to respect convective heating limits without disrupting the rapid Jacobian convergence. Finally, the integration architecture handles a continuous physics "idle" during the Karman line transition, propagating a $J_2$-perturbed Keplerian orbital coast before re-engaging quaternion kinematics for the terminal hypersonic reentry phase—all within a single, unbroken, GPU-accelerated numerical integration run.
\end{abstract}

\pagebreak

\section{Introduction}
\label{sec:intro}
The modeling and simulation of Hypersonic Glide Vehicles (HGVs) and Fractional Orbital Bombardment Systems (FOBS) present a highly coupled, multi-physics challenge that fundamentally strains traditional computational aerospace methods. Legacy trajectory and aeroelastic analysis often relies on discrete Finite Element Methods (FEM) and rigidly decoupled aerodynamic look-up tables. While effective for localized, static structural validation, these discrete methodologies introduce discontinuous boundary conditions, require mid-flight matrix inversions, and mandate prohibitively long computational wall-times. This renders them highly sub-optimal for the rapid, large-scale Monte-Carlo dispersion studies required in modern trajectory design and autonomous guidance validation.

This manuscript introduces PASSES-HGV, a comprehensive expansion of the original spectral multi-physics framework. By elevating the core fixed-grid, quasi-static architecture to a full 6-Degree-of-Freedom (6-DOF) kinematic environment, this framework establishes a continuous mathematical pipeline capable of seamlessly simulating lifting-body aerodynamics and exo-atmospheric orbital mechanics. To circumvent the computational bottlenecks of traditional meshing, the HGV structural and thermodynamic domains are mapped onto a 2D bivariate Chebyshev tensor product space. The complex, direction-dependent torsional and transverse bending modes inherent to waverider geometries are governed continuously by the Knox-Xi tensor, ensuring an exact, $C^\infty$ differentiable representation of anisotropic flexural rigidity across the entire vehicle planform.

Furthermore, this framework integrates a fully autonomous, unclassified hypersonic aerothermodynamic model. By coupling modified Newtonian flow approximations and Sutton-Graves convective heating directly to the structural state vector, the right-hand side forcing operators remain perfectly continuous and suitable for massive GPU acceleration. Navigational control is achieved via an embedded numerical Predictor-Corrector guidance algorithm utilizing continuous exterior penalty functions, allowing the vehicle to actively manage kinetic energy while autonomously respecting strict thermal and skip-out constraints. Ultimately, PASSES-HGV demonstrates the capability to execute an unbroken, highly parallelized numerical integration spanning from a dynamic atmospheric boost phase, through a $J_2$-perturbed fractional Keplerian orbit, and deep into the extreme aerothermal regime of hypersonic reentry.

\section{The Expanded Global State Vector}
\label{sec:global_state_vector}

\subsection{The 6-DOF Attitude State}
To prevent singularity issues (gimbal lock) when the HGV executes aggressive, near-vertical dive maneuvers or exo-atmospheric tumbles, we parameterize the vehicle's attitude using a unit quaternion $\mathbf{q} = [q_0, q_1, q_2, q_3]^T$, where $q_0$ is the scalar component.

The time evolution of the attitude is coupled to the body-frame angular velocity $\boldsymbol{\omega}_B = [p, q, r]^T$ via the quaternion kinematic differential equation:

$$\dot{\mathbf{q}} = \frac{1}{2} \mathbf{Q}(\boldsymbol{\omega}_B) \mathbf{q}$$

where $\mathbf{Q}(\boldsymbol{\omega}_B)$ is the skew-symmetric matrix multiplier for the angular rates. This differential equation is integrated concurrently with the structural state $\mathbf{w}$ within the RKF45 solver.

\subsection{The Direction Cosine Matrix (DCM)}
To calculate aerodynamic forces, we must transform the Earth-relative freestream velocity vector, $\mathbf{v}_\infty$, from the inertial frame (or Earth-Centered Earth-Fixed frame) into the vehicle's body frame.

We map the quaternion $\mathbf{q}$ to the standard rotation matrix, $\mathbf{T}_I^B(\mathbf{q})$:

$$\mathbf{T}_I^B(\mathbf{q}) = \begin{bmatrix} q_0^2 + q_1^2 - q_2^2 - q_3^2 & 2(q_1 q_2 + q_0 q_3) & 2(q_1 q_3 - q_0 q_2) \\ 2(q_1 q_2 - q_0 q_3) & q_0^2 - q_1^2 + q_2^2 - q_3^2 & 2(q_2 q_3 + q_0 q_1) \\ 2(q_1 q_3 + q_0 q_2) & 2(q_2 q_3 - q_0 q_1) & q_0^2 - q_1^2 - q_2^2 + q_3^2 \end{bmatrix}$$

The body-frame velocity vector is then simply $\mathbf{v}_B = \mathbf{T}_I^B \mathbf{v}_\infty$. From $\mathbf{v}_B$, the global angle of attack ($\alpha$) and sideslip angle ($\beta$) are instantaneously derived.

\subsection{Local Incidence Angle Mapping}
Because the waverider is not a flat plate—it has complex curvature and continuous aeroelastic deformation—the global angle of attack $\alpha$ is insufficient. We need the local incidence angle, $\delta_c(\xi, \eta)$, at every node on the bivariate Chebyshev grid.

For a rigid, undeflected shape, the local surface normal vector $\mathbf{n}_0(\xi, \eta)$ is known from the CAD transcription. However, under high dynamic pressure, the spectral plate bends by $\mathbf{w}(\xi, \eta)$. The *deformed* local surface normal $\mathbf{n}(\xi, \eta, t)$ is a function of the spatial derivatives of the structural state (which we have readily available via our Kronecker matrices $\mathbf{D}_x$ and $\mathbf{D}_y$).

The local incidence angle is thus the dot product of the normalized body-frame velocity vector and the deformed surface normal:

$$\sin(\delta_c(\xi, \eta)) = \mathbf{n}(\xi, \eta, t) \cdot \frac{\mathbf{v}_B}{\vert{}\mathbf{v}_B\vert{}}$$

\subsection{Constructing the RHS Forcing Vector}
With $\delta_c$ defined at every grid point, we can instantly map the modified Newtonian pressure field to the vehicle surface. Any node located in the aerodynamic shadow (where $\delta_c < 0$) is assigned zero pressure. For the windward nodes, the local pressure coefficient is:

$$C_p(\xi, \eta) = C_{p,max} \sin^2(\delta_c(\xi, \eta))$$

The continuous dynamic pressure $q_{dyn} = \frac{1}{2} \rho \vert{}\mathbf{v}_\infty\vert{}^2$ scales this field into physical dimensional pressure. Flattening this 2D continuous pressure field over the bivariate grid yields our highly sought-after quasi-static RHS forcing vector:

$$\mathbf{Q}_{aero} = q_{dyn} \mathbf{C}_p$$

This couples the rigid-body 6-DOF kinematics, the freestream environment, and the flexible structural dynamics into one seamless, continuous mathematical pipeline.

\section{High-Fidelity Hypersonic Aerothermodynamics}
\label{sec:hypersonic_aerothermodynamics}
To design quantitative Thermal Protection Systems (TPS) for hypersonic waveriders, engineering-level correlations are insufficient, as they fail to capture finite-rate chemistry, molecular dissociation, and radiative heat transfer present at extreme flight regimes. However, embedding a full Navier-Stokes Computational Fluid Dynamics (CFD) solver within an RKF45 integration loop breaks the computational throughput required for real-time guidance projections. We resolve this by substituting high-fidelity analytical correlations that mimic real-gas thermochemistry directly into the continuous right-hand side (RHS) forcing vector.

\subsection{Real-Gas Convective Heating: The Fay-Riddell Formulation}
We anchor the continuous thermal model using the Fay-Riddell formulation for the stagnation point convective heat flux. This accounts for boundary layer equilibrium chemistry and surface catalytic recombination (the exothermic release of energy when dissociated atmospheric oxygen and nitrogen recombine on the vehicle's leading edges).

The stagnation point convective heat flux is governed by the enthalpy-driven relation:

$$\dot{q}_{s, conv} = 0.763 \text{Pr}^{-0.6} (\rho_e \mu_e)^{0.4} (\rho_w \mu_w)^{0.1} \sqrt{\left(\frac{du_e}{dx}\right)_s} (h_{0e} - h_w) \left[ 1 + (Le^\beta - 1)\frac{h_D}{h_{0e}} \right]$$

where $\rho_e, \mu_e$ are the density and viscosity at the boundary layer edge, $\rho_w, \mu_w$ are the properties at the vehicle wall, $h_{0e}$ is the total stagnation enthalpy, $h_w$ is the wall enthalpy, and $h_D$ is the dissociation enthalpy. The Lewis number ($Le$) term explicitly captures the diffusion of atomic species toward the TPS boundary.

\subsection{Radiative Heat Transfer: The Tauber-Sutton Formulation}
At velocities exceeding 5 km/s, the shock layer compresses the atmosphere to temperatures where thermal radiation becomes a dominant, and sometimes overriding, heating mechanism. Unlike convective heating, which scales proportionally to $V^3$, radiative heating scales roughly with $V^8$ to $V^{10}$ and is strongly volume-dependent.

We model the stagnation radiative heat flux continuously using the Tauber-Sutton empirical formulation:

$$\dot{q}_{s, rad} = C \left( \frac{\rho_\infty}{\rho_0} \right)^n f(V_\infty) R_{eff}^a$$

where $C, n,$ and $a$ are atmospheric curve-fit constants, and $f(V_\infty)$ is a high-order velocity polynomial. The effective nose radius $R_{eff}$ is dynamically extracted via the 2D CAD transcription's local curvature at the stagnation node, allowing $\dot{q}_{s, rad}$ to adapt instantaneously as the ablation matrix changes the vehicle geometry.

\subsection{The Continuous 2D Total Heat Flux Tensor}
The total stagnation heating $\dot{q}_s = \dot{q}_{s, conv} + \dot{q}_{s, rad}$ is then mapped across the entire bivariate grid using the incidence angle $\delta_c(\xi, \eta)$ computed from the 6-DOF quaternion kinematics.

For the windward nodes ($\delta_c > 0$), a continuous modified Lees' distribution governs the local heat flux:

$$\dot{q}_{total}(\xi, \eta) = \dot{q}_s \sin(\delta_c(\xi, \eta)) \left( \frac{R_{eff}}{x(\xi)} \right)^{0.5}$$

This generates a massive 2D matrix of heat flux values across the waverider's surface: $\mathbf{Q}_{thermal} = \text{diag}(\dot{q}_{total}(\xi_i, \eta_j))$. Because this thermal RHS is formulated strictly through continuous tensor-product operations, it maintains $C^\infty$ continuity and maps flawlessly to parallelized matrix operations in C/C++. The total heat flux acts as the time-varying Neumann boundary condition for the smoothed enthalpy ablation matrix, dynamically coupling the real-gas thermochemistry to the structural recession of the vehicle.

\section{Structural Generality: Mindlin-Reissner Anisotropic Plate Theory}
\label{sec:structural_generality}
Transitioning from a slender rocket to a hypersonic lifting body means the vehicle is no longer governed by simple 1D beam mechanics. Furthermore, because waveriders exhibit highly variable thickness profiles—thick at the root and thin at the leading edges—classical Kirchhoff-Love thin-plate theory is insufficient. It explicitly ignores transverse shear deformations, which become the dominant mode of aeroelastic flutter in the thicker sections of the hull under high aerodynamic loading.

To resolve this, we elevate the structural framework to a generalized Mindlin-Reissner anisotropic plate formulation. This allows the cross-section to undergo shear deformation independently of the transverse bending, capturing the true 3D aeroelastic response of the vehicle without resorting to the computational bloat of discrete 3D solid finite elements.

\subsection{The Expanded 3-DOF Structural State}
In the Mindlin-Reissner framework, the continuous structural state at any spatial coordinate $(\xi, \eta)$ is no longer a scalar deflection $w$. It expands into a $3 \times 1$ vector:

$$\mathbf{u}(\xi, \eta, t) = \begin{bmatrix} w \\ \phi_x \\ \phi_y \end{bmatrix}$$

where $w$ is the transverse deflection, and $\phi_x, \phi_y$ are the independent normal rotations about the $y$ and $x$ axes, respectively.

Because the rotations are no longer strictly locked to the gradient of the surface, the localized transverse shear strains are defined continuously across the domain as the difference between the geometric slope and the structural rotation:

$$\gamma_{xz} = \phi_x + \frac{\partial w}{\partial x}, \quad \gamma_{yz} = \phi_y + \frac{\partial w}{\partial y}$$

\subsection{The Block-Kronecker Stiffness Matrix}
We map the vehicle's planform onto a canonical 2D computational domain using the bivariate Gauss-Lobatto grid. By formulating the Mindlin-Reissner governing equations as a system of coupled, block-sparse Kronecker operators, the structural domain is cast as a pure applied mathematics framework, bypassing the discontinuous assembly steps of traditional finite element methods.

The flexural rigidity tensor $\mathbf{D}(\xi, \eta)$ is now augmented by the transverse shear rigidity tensor $\mathbf{A}(\xi, \eta)$, where $A_{11} = \kappa^2 G_{xz} h$ and $A_{22} = \kappa^2 G_{yz} h$, with $\kappa^2 = 5/6$ acting as the shear correction factor.

When flattened over the $(N_x+1) \times (N_y+1)$ grid, the total structural state vector $\mathbf{U}$ expands to a size of $3N$. The global stiffness matrix $\mathbf{K}$ becomes a massive $3N \times 3N$ block matrix built entirely from tensor products of 1D Chebyshev differentiation matrices ($\mathbf{D}_x = \mathbf{D}_{N_x} \otimes \mathbf{I}_{N_y}$):

$$\mathbf{K} = \begin{bmatrix} \mathbf{K}_{ww} & \mathbf{K}_{w\phi_x} & \mathbf{K}_{w\phi_y} \\ \mathbf{K}_{\phi_x w} & \mathbf{K}_{\phi_x \phi_x} & \mathbf{K}_{\phi_x \phi_y} \\ \mathbf{K}_{\phi_y w} & \mathbf{K}_{\phi_y \phi_x} & \mathbf{K}_{\phi_y \phi_y} \end{bmatrix}$$

Every block is analytically defined. For example, the shear-coupling block governing the interaction between deflection and $x$-rotation is constructed as:

$$\mathbf{K}_{w\phi_x} = \mathbf{D}_x^T \mathbf{A}_{11, flex}$$

\subsection{GPU-Accelerated Dynamic Propagation}
The global ODE system for the structural domain remains mathematically identical in form to a standard linear oscillator:

$$\mathbf{M} \mathbf{\ddot{U}} + \mathbf{K} \mathbf{U} = \mathbf{Q}_{aero}$$

However, the mass matrix $\mathbf{M}$ is now a block-diagonal operator incorporating both the translational mass $m(\xi, \eta)$ and the rotary inertias $I(\xi, \eta)$ of the thick plate.

Because the global stiffness matrix $\mathbf{K}$ is pre-computed as a static block tensor, the RKF45 integration maps flawlessly to CUDA-accelerated sparse matrix-vector multiplications. This bypasses the thread-divergence bottlenecks often encountered in highly coupled physics simulations or custom parallelized raytracers, ensuring the structural dynamics solve fast enough to support the embedded flight computer's real-time shadow predictions.

\section{Boundary Conditions}
In the 1D model, we embedded the free-free constraints (zero moment, zero shear) directly into the stiffness matrix $\mathbf{K}$ via explicit row replacement. Expanding this to a 2D lifting body is algebraically intense because a free-floating waverider is governed by the Kirchhoff-Love boundary conditions along its entire perimeter, which couples the longitudinal and transverse bending modes.

\subsection{Anisotropic Rigidity Tensor}

Because the hypersonic glide vehicle relies on a blended, multi-material aerodynamic shell, its resistance to bending and torsion is highly direction-dependent. Rather than relying on simplified isotropic constants, we define the vehicle's spatially continuous anisotropic flexural rigidity tensor, $\mathbf{D}(\xi, \eta)$.

This tensor maps the continuous planform to four primary stiffness distributions:

\begin{itemize}
    \item $\mathbf{D}_{11}(\xi, \eta)$: Longitudinal bending stiffness
    \item $\mathbf{D}_{22}(\xi, \eta)$: Transverse (chordwise) bending stiffness
    \item $\mathbf{D}_{12}(\xi, \eta)$: Poisson-effect coupling stiffness
    \item $\mathbf{D}_{66}(\xi, \eta)$: Torsional (twisting) stiffness
\end{itemize}

When flattened into diagonalized matrices over the bivariate Gauss-Lobatto grid (e.g., $\mathbf{D}_{11} = \text{diag}(\mathbf{D}_{11}(\xi_i, \eta_j))$), these components become the core weighting matrices nestled between the Kronecker differentiation operators to form the unconstrained global stiffness matrix $\mathbf{K}_{unc}$.

\subsection{The Free-Free Kirchhoff Boundary Conditions}
For an unconstrained plate floating in the atmosphere, the edges must experience zero bending moment and zero effective shear. Let's isolate the boundaries along the longitudinal edges ($x = \pm L_x$, corresponding to canonical nodes $\xi = \pm 1$).

The zero bending moment ($M_x = 0$) constraint is coupled by the Knox-Xi tensor:

$$M_x = -\mathbf{D}_{11} \frac{\partial^2 w}{\partial x^2} - \mathbf{D}_{12} \frac{\partial^2 w}{\partial y^2} = 0$$

The zero effective Kirchhoff shear ($V_x = 0$) constraint requires the sum of the transverse shear and the gradient of the twisting moment to vanish:

$$V_x = -\frac{\partial}{\partial x} \left( \mathbf{D}_{11} \frac{\partial^2 w}{\partial x^2} + \mathbf{D}_{12} \frac{\partial^2 w}{\partial y^2} \right) - 2 \frac{\partial}{\partial y} \left( \mathbf{D}_{66} \frac{\partial^2 w}{\partial x \partial y} \right) = 0$$

\subsection{Spectral Row Replacement in the Flattened Tensor Space}
To enforce these constraints without relying on discrete ghost nodes, we must embed them directly into the quasi-static matrices.

In a 1D grid, the boundaries are simply indices $0$ and $N$. In the 2D tensor product grid of size $(N_x+1) \times (N_y+1)$, the structural state vector $\mathbf{w}$ is flattened. We define a boundary index set $\mathcal{B}$ containing the 1D indices of all nodes lying on the perimeter (the top/bottom rows and left/right columns of the logical grid).

For every index $b \in \mathcal{B}$:
\begin{itemize}
    \item \textbf{Decouple Inertia:} The $b$-th row of the global mass matrix is overwritten with zeros ($\mathbf{M}_{b, :} = \mathbf{0}$), and the corresponding row in the forcing vector is zeroed ($\mathbf{Q}_b = 0$).
    \item \textbf{Overwrite Stiffness:} The $b$-th row of the global stiffness matrix $\mathbf{K}$ is replaced by the discrete spectral boundary operators.
\end{itemize}

For a node on the edge $\xi = 1$, the moment row replacement using the pre-computed Kronecker matrices becomes:

$$\mathbf{K}_{b, :} = \left[ \mathbf{D}_{11} \mathbf{D}_x^2 + \mathbf{D}_{12} \mathbf{D}_y^2 \right]_{b, :}$$

And the shear row replacement for the adjacent inner boundary layer becomes:

$$\mathbf{K}_{b-1, :} = \left[ \mathbf{D}_x \left( \mathbf{D}_{11} \mathbf{D}_x^2 + \mathbf{D}_{12} \mathbf{D}_y^2 \right) + 2 \mathbf{D}_y \left( \mathbf{D}_{66} \mathbf{D}_{xy} \right) \right]_{b-1, :}$$

By sweeping through the perimeter set $\mathcal{B}$ and injecting these exact Kronecker rows into $\mathbf{K}$, the plate's free-free physics are mathematically locked into the RHS operator before the flight simulation even begins. The RKF45 solver will naturally propagate the twisting and bending waves to the edges and perfectly reflect them back into the vehicle structure, avoiding the Gibbs phenomenon.

\section{GNC: Energy Management vs. Terminal Intercept}
\label{sec:gnc}

\subsection{The HGV Flight Corridor Constraints}
Before predicting the trajectory, the algorithm must understand its bounds. An HGV must thread the needle between two physical limits, defining the "Equilibrium Glide Phase":

\begin{itemize}
    \item \textbf{Upper Bound (Thermal Limit):} The vehicle cannot dive too steeply, or the Sutton-Graves convective heat flux $\dot{q}_{conv}$ and dynamic pressure $q_{dyn}$ will exceed structural limits.
    \item \textbf{Lower Bound (Skip-Out Limit):} The vehicle cannot fly too high, or the aerodynamic lift will exceed the apparent weight, causing the vehicle to skip back out of the atmosphere.
\end{itemize}

The nominal lift required to maintain equilibrium flight is:

$$L_{eq} = m \left( g - \frac{V^2}{r} \right)$$

where $V$ is Earth-relative velocity and $r$ is the radial distance from the Earth's center. The bank angle $\sigma$ directs this lift vector. The vertical component keeps the vehicle in the corridor, while the horizontal component steers it laterally.

\subsection{The Numerical Predictor (The "Shadow" Solver)}

Traditional ballistic missiles use analytical approximations (like Keplerian ellipses) to predict the impact point. HGVs cannot do this because aerodynamic forces dominate over gravity.

Instead of relying on analytical approximations, the flight computer can spawn a lightweight, "shadow" version of our RKF45 solver in real-time. Because our architecture is designed around C/C++ and CUDA acceleration, these forward-projections can be executed in milliseconds. We simply dedicate a subset of GPU threads to integrate a 3-DOF point-mass version of the equations of motion from the current time $t_0$ to the terminal target velocity $V_f$.

During this forward projection, the shadow solver integrates:

$$\mathbf{\dot{r}} = \mathbf{v}$$

$$\mathbf{\dot{v}} = \mathbf{g}(\mathbf{r}) + \frac{1}{m} \mathbf{T}_B^I \mathbf{F}_{aero}(\alpha, \sigma)$$

using a parameterized guess for the bank angle profile, often defined as a linear function of specific energy $E$:

$$\sigma(E) = \sigma_0 + k (E - E_0)$$

\subsection{Successive Convexification and the Optimal Control Problem}
While numerical Predictor-Corrector algorithms utilizing Secant or Newton-Raphson updates are computationally lightweight, they cannot rigorously enforce hard boundary constraints. To prevent constraint violation during aggressive maneuvering, we reformulate the trajectory generation as a constrained Optimal Control Problem (OCP).

The objective is to find the continuous control history—specifically the bank angle profile $\sigma(t)$ and angle of attack $\alpha(t)$—that minimizes a cost function $J$ (e.g., maximizing terminal velocity $V_f$ or minimizing integrated heat load) subject to the highly non-linear equations of motion and strict flight corridor boundaries.

We define the 6-DOF kinematic and structural state vector as $\mathbf{x}(t)$ and the control vector as $\mathbf{u}(t) = [\alpha(t), \sigma(t)]^T$. The continuous non-linear OCP is formulated as:

$$\min_{\mathbf{u}(t)} \quad J(\mathbf{x}(t_f), \mathbf{u}(t))$$

$$\text{subject to:} \quad \dot{\mathbf{x}}(t) = f(\mathbf{x}(t), \mathbf{u}(t))$$

$$\mathbf{x}(t_0) = \mathbf{x}_0, \quad \psi(\mathbf{x}(t_f)) = \mathbf{0}$$

$$g(\mathbf{x}(t), \mathbf{u}(t)) \le 0$$

where $\psi(\mathbf{x}(t_f))$ represents the terminal intercept constraints (downrange and crossrange targets), and $g(\mathbf{x}, \mathbf{u})$ encapsulates the continuous state-control inequality constraints.

\subsection{Non-Linear Constraint Formulation}
The primary boundaries governing the hypersonic glide phase are structural and thermodynamic. We explicitly define these constraints as functions of the state vector:

\begin{enumerate}
    \item \textbf{Thermal Heating Limit:} The Fay-Riddell stagnation heat flux must not exceed the structural limit of the TPS: $$\dot{q}_{s, conv}(\mathbf{x}) - \dot{q}_{max} \le 0$$
    \item \textbf{Dynamic Pressure Limit:} To prevent aeroelastic failure, the dynamic pressure must remain bounded: $$\frac{1}{2} \rho(\mathbf{x}) V(\mathbf{x})^2 - q_{max} \le 0$$
    \item \textbf{Control Effort Bounds:} The aerodynamic control surfaces have physical actuation limits: $$\vert{}\sigma(t)\vert{} - \sigma_{max} \le 0$$
\end{enumerate}

Because the dynamics $f(\mathbf{x}, \mathbf{u})$ and constraints $g(\mathbf{x}, \mathbf{u})$ are highly non-linear and non-convex, a standard optimization algorithm cannot guarantee a global minimum and will likely fail to converge within the millisecond timeframe required for active flight control.

\subsection{Linearization and Second-Order Cone Programming (SOCP)}
To resolve this, we employ Successive Convexification. The algorithm linearizes the continuous non-linear dynamics around a known reference trajectory, $(\mathbf{x}^{(k)}(t), \mathbf{u}^{(k)}(t))$, using a first-order Taylor series expansion:

$$\dot{\mathbf{x}}(t) \approx f(\mathbf{x}^{(k)}, \mathbf{u}^{(k)}) + \mathbf{A}(t) \delta\mathbf{x} + \mathbf{B}(t) \delta\mathbf{u}$$

where $\mathbf{A}(t) = \frac{\partial f}{\partial \mathbf{x}}$ and $\mathbf{B}(t) = \frac{\partial f}{\partial \mathbf{u}}$ are the continuous-time Jacobian matrices evaluated along the reference trajectory, and $\delta\mathbf{x} = \mathbf{x} - \mathbf{x}^{(k)}$.

We discretize this linearized system over $N$ temporal nodes using Zero-Order Hold (ZOH) or Runge-Kutta integration, transforming the dynamics into a set of linear equality constraints:

$$\mathbf{x}_{i+1} = \mathbf{A}_i \mathbf{x}_i + \mathbf{B}_i \mathbf{u}_i + \mathbf{z}_i, \quad i = 1, \dots, N-1$$

The non-linear state constraints (heating, dynamic pressure) are identically linearized into convex half-spaces. To prevent the optimization from commanding a trajectory that wildly diverges from the reference path (where the linearization is no longer valid), we introduce a convex Trust Region penalty bounding the step size:

$$\Vert \mathbf{x} - \mathbf{x}^{(k)} \Vert_2 \le \Delta_{trust}$$

\subsection{The Autonomous SC Loop}
The entire OCP is now cast as a Second-Order Cone Program (SOCP), a subclass of convex optimization that can be solved deterministically with zero risk of local minima. The autonomous guidance architecture executes the following loop:

\begin{enumerate}
    \item \textbf{Initialize:} Generate a kinematically feasible (but suboptimal) reference trajectory using a rapid 3-DOF shadow solver.
    \item \textbf{Convexify:} Linearize the dynamics and state constraints around the current trajectory.
    \item \textbf{Optimize:} Execute an onboard interior-point solver (e.g., ECOS or custom CUDA-accelerated primal-dual solver) to find the optimal control $\mathbf{u}^{(k+1)}$ and state $\mathbf{x}^{(k+1)}$ for the convex subproblem.
    \item \textbf{Update:} Set the new trajectory as the reference. If the change in the cost function $\Vert J^{(k+1)} - J^{(k)} \Vert$ is below a defined tolerance, the algorithm has converged.
    \item \textbf{Actuate:} Command the first control input $\mathbf{u}_1$ to the 6-DOF quaternion controller.
\end{enumerate}

By embedding this Successive Convexification algorithm within the fixed-grid spectral framework, the waverider autonomously rides the very edge of the thermal boundary, maximizing kinetic energy retention while mathematically guaranteeing it will never violate the structural heating limits of the carbon-carbon leading edges.

\subsection{Plasma Sheath Ionization and RF Attenuation}

During atmospheric entry at Mach numbers exceeding 15, the shock layer compresses and heats the ambient air to temperatures capable of molecular dissociation and atomic ionization. This creates a plasma sheath enveloping the vehicle, saturated with free electrons.

To determine when telemetry blackout occurs, we dynamically compute the local electron number density $n_e$ at the boundary layer edge by coupling the stagnation enthalpy states from the Fay-Riddell correlations to the Saha ionization equation. The critical plasma frequency $\omega_p$, which dictates the electromagnetic opacity of the shock layer, is continuously evaluated as:

$$\omega_p = \sqrt{\frac{n_e e^2}{m_e \epsilon_0}}$$

where $e$ is the elementary charge, $m_e$ is the electron mass, and $\epsilon_0$ is the vacuum permittivity.

Whenever the plasma frequency exceeds the transmission frequency of the vehicle's telemetry network (e.g., L-band GNSS at $\approx 1.5$ GHz), the incident electromagnetic waves are completely reflected. The framework models this mathematically by driving the signal-to-noise ratio to zero, completely severing external navigation updates and initiating a prolonged RF blackout phase.

\subsection{The Ionization-Aware Extended Kalman Filter (EKF)}
To maintain state estimation during blackout, the flight computer relies on a tightly-coupled, ionization-aware Extended Kalman Filter (EKF). The navigation state vector $\mathbf{x}_{nav} \in \R^{15}$ is formulated to track kinematics and sensor drift:

$$\mathbf{x}_{nav} = \begin{bmatrix} \mathbf{r} \\ \mathbf{v} \\ \mathbf{q} \\ \mathbf{b}_a \\ \mathbf{b}_g \end{bmatrix}$$

where $\mathbf{r}, \mathbf{v}$ are the ECEF position and velocity vectors, $\mathbf{q}$ is the attitude quaternion, and $\mathbf{b}_a, \mathbf{b}_g$ are the continuous drift biases of the onboard Inertial Measurement Unit (IMU) accelerometers and gyroscopes.

The EKF prediction step integrates the continuous IMU dynamics:

$$\dot{\hat{\mathbf{x}}}_{nav}(t) = f_{nav}(\hat{\mathbf{x}}_{nav}(t), \mathbf{u}_{imu}(t))$$

Simultaneously, the state error covariance matrix $\mathbf{P}(t)$ is propagated via the continuous-time Riccati equation:

$$\dot{\mathbf{P}}(t) = \mathbf{F}(t) \mathbf{P}(t) + \mathbf{P}(t) \mathbf{F}(t)^T + \mathbf{Q}_{noise}$$

where $\mathbf{F}(t)$ is the navigation Jacobian and $\mathbf{Q}_{noise}$ represents the stochastic noise profile of the inertial sensors.

\subsection{Dynamic Measurement Gating}
The true innovation of the PASSES-HGV navigation framework is its dynamic measurement gating. The discrete measurement update equation is defined as:

$$\mathbf{K}_k = \mathbf{P}_k^- \mathbf{H}_k^T (\mathbf{H}_k \mathbf{P}_k^- \mathbf{H}_k^T + \mathbf{R}_k)^{-1}$$

$$\hat{\mathbf{x}}_k^+ = \hat{\mathbf{x}}_k^- + \mathbf{K}_k (\mathbf{z}_k - h(\hat{\mathbf{x}}_k^-))$$

$$\mathbf{P}_k^+ = (\mathbf{I} - \mathbf{K}_k \mathbf{H}_k) \mathbf{P}_k^-$$

When the vehicle operates in standard atmospheric regimes ($\omega_p < \omega_{GNSS}$), the observation matrix $\mathbf{H}_k$ actively maps external GNSS pseudoranges into the state update, clamping the covariance $\mathbf{P}$.

However, when the thermochemical module flags an active plasma blackout ($\omega_p \ge \omega_{GNSS}$), the EKF autonomously zeroes the corresponding rows in the observation matrix $\mathbf{H}_k \to \mathbf{0}$. The Kalman gain $\mathbf{K}_k$ drops to zero for external sensors, and the EKF transitions into a pure inertial coasting mode. During this blackout window, the state covariance $\mathbf{P}(t)$ balloons quadratically due to the unobservable IMU biases $\mathbf{b}_a$ and $\mathbf{b}_g$.

This directly feeds back into the Successive Convexification algorithm (Section 6.5). If the projected covariance footprint $\mathbf{P}(t_f)$ exceeds the target intercept radius, the optimal control solver will autonomously command a "pull-up" maneuver. This bleeds off velocity to lower the stagnation temperature, extinguishing the plasma sheath and forcing a GNSS re-acquisition window, all while maintaining the mathematically guaranteed thermal constraints of the airframe.

\section{Statistical Validation and GPU-Accelerated Monte Carlo Dispersion}
\label{sec:monte_carlo_validation}

While deterministic trajectory propagation establishes the mathematical consistency of the continuous multi-physics operators, operational hypersonic vehicles operate in highly stochastic environments. Variations in atmospheric density, aerodynamic coefficients, and mass properties induce significant terminal dispersion. To transition the framework from a theoretical solver to a production-level validation architecture, we formulate a parallelized Monte Carlo dispersion pipeline.

\subsection{The Stochastic Parameter Space}

We define a stochastic parameter vector $\mathbf{p} \in \R^k$ that encapsulates the primary environmental and vehicle uncertainties:

$$\mathbf{p} = \begin{bmatrix} \delta \rho_\infty \\ \delta C_L \\ \delta C_D \\ \delta m \\ \delta J_2 \end{bmatrix}$$

where $\delta$ represents the perturbation from the nominal scalar value. We assume the parameter space is governed by a multivariate Gaussian distribution with mean $\boldsymbol{\mu}_p = \mathbf{0}$ and a known, dense covariance matrix $\boldsymbol{\Sigma}_p$, which accounts for cross-correlation between aerodynamic properties.

To generate the Monte Carlo test matrix, we sample $N_{sim}$ discrete realizations of the parameter vector via the Cholesky decomposition of the covariance matrix:

$$\mathbf{p}^{(i)} = \boldsymbol{\mu}_p + \mathbf{L} \mathbf{z}^{(i)}, \quad i = 1, 2, \dots, N_{sim}$$

where $\boldsymbol{\Sigma}_p = \mathbf{L} \mathbf{L}^T$ and $\mathbf{z}^{(i)} \sim \mathcal{N}(\mathbf{0}, \mathbf{I})$.

\subsection{Batched-Tensor CUDA Architecture}

Traditional Monte Carlo validation relies on sequential `for`-loops over the numerical integrator, which is computationally prohibitive for high-fidelity aeroelastic models. Because the PASSES-HGV framework is built exclusively on static Kronecker product matrices ($\mathbf{M}$ and $\mathbf{K}$), we can vector-ize the entire simulation suite.

We expand the 6-DOF structural and kinematic state vector $\mathbf{X}$ into a rank-3 state tensor $\boldsymbol{\mathcal{X}}(t) \in \R^{N_{sim} \times N_{state}}$. The RKF45 integration loop is executed exactly once, but the underlying matrix multiplications are cast as batched-tensor operations.

At each time step, a highly parallelized CUDA kernel assigns a dedicated thread block to each of the $N_{sim}$ trajectories. The continuous RHS aerodynamic forcing tensor $\boldsymbol{\mathcal{Q}}_{aero}(\boldsymbol{\mathcal{X}}, \mathbf{p})$ is evaluated simultaneously for all dispersed states, bypassing thread divergence and allowing the computation of 10,000+ stochastic hypersonic glides in the time a standard CPU solver requires for a single deterministic run.

\subsection{Terminal Footprint and Covariance Extraction}
The output of the batched integration is a dispersion cloud of terminal impact coordinates (downrange $R$ and crossrange $C$). We extract the statistical footprint by calculating the sample mean impact point:

$$\boldsymbol{\mu}_{impact} = \frac{1}{N_{sim}} \sum_{i=1}^{N_{sim}} \begin{bmatrix} R_{pred}^{(i)} \\ C_{pred}^{(i)} \end{bmatrix}$$

and the $2 \times 2$ terminal covariance matrix:

$$\boldsymbol{\Sigma}_{impact} = \frac{1}{N_{sim} - 1} \sum_{i=1}^{N_{sim}} \left( \begin{bmatrix} R_{pred}^{(i)} \\ C_{pred}^{(i)} \end{bmatrix} - \boldsymbol{\mu}_{impact} \right) \left( \begin{bmatrix} R_{pred}^{(i)} \\ C_{pred}^{(i)} \end{bmatrix} - \boldsymbol{\mu}_{impact} \right)^T$$

This terminal covariance matrix maps directly to the 3-sigma equiprobability error ellipse on the planetary surface. By embedding this statistical architecture natively within the fixed-grid spectral solver, the framework provides the exact probabilistic validation products required to benchmark against high-fidelity, coupled CFD/CSD computational codes.

\section{Orbital Mechanics and FOBS}
\label{sec:orbital_mechanics_fobs}

\subsection{The Continuous Atmospheric Fade (The "Physics Idle")}
As the vehicle ascends past the Karman line (approx. $100$ km), the atmospheric density $\rho(h)$ exponentially decays toward zero. Because our entire right-hand side (RHS) forcing vector is scaled by the dynamic pressure $q_{dyn} = \frac{1}{2} \rho \vert{}\mathbf{v}_\infty\vert{}^2$, the aerodynamic pressure $\mathbf{Q}_{aero}$ naturally and continuously vanishes.

Simultaneously, the Sutton-Graves convective heat flux $\dot{q}_{conv}$, which relies on $\sqrt{\rho}$, decays to zero.

Mathematically, the multi-physics PDE system gracefully degrades into an unforced system:

$$\mathbf{M} \mathbf{\ddot{w}} + \mathbf{K} \mathbf{w} = \mathbf{0}$$

With no external forcing, internal structural damping quickly dissipates any residual aeroelastic vibrations. The structural velocity $\mathbf{\dot{w}}$ approaches zero, and the vehicle enters a rigid-body state. Because the structural derivatives are no longer exhibiting high-frequency oscillations, the RKF45 time-stepping algorithm will automatically detect the drop in stiffness and massively expand its integration step size ($\Delta t$), allowing the flight computer to simulate a 45-minute orbital coast in a fraction of a second.

\subsection{$J_2$ Perturbed Orbital Mechanics}
In a true FOBS trajectory, a simple spherical Earth gravity model ($\mathbf{g} = -\frac{\mu}{r^3} \mathbf{r}$) is grossly insufficient. A low Earth orbit (LEO) coast phase is highly sensitive to the Earth's oblateness (the equatorial bulge). Over a fractional orbit, this induces secular variations in the trajectory that will cause catastrophic crossrange errors during reentry if ignored.

We must update the translational kinematics to include the $J_2$ geopotential harmonic. In the Earth-Centered Inertial (ECI) Cartesian frame, the gravitational acceleration vector $\mathbf{g}(\mathbf{r})$ becomes:

$$\mathbf{g}(\mathbf{r}) = -\frac{\mu}{\vert{}\mathbf{r}\vert{}^3} \mathbf{r} + \mathbf{g}_{J_2}(\mathbf{r})$$

The $J_2$ perturbation vector is derived from the gradient of the geopotential function:

$$\mathbf{g}_{J_2}(\mathbf{r}) = -\frac{3}{2} J_2 \left(\frac{\mu}{\vert{}\mathbf{r}\vert{}^2}\right) \left(\frac{R_E}{\vert{}\mathbf{r}\vert{}}\right)^2 \begin{bmatrix} \left(1 - 5\frac{z^2}{\vert{}\mathbf{r}\vert{}^2}\right) \frac{x}{\vert{}\mathbf{r}\vert{}} \\ \left(1 - 5\frac{z^2}{\vert{}\mathbf{r}\vert{}^2}\right) \frac{y}{\vert{}\mathbf{r}\vert{}} \\ \left(3 - 5\frac{z^2}{\vert{}\mathbf{r}\vert{}^2}\right) \frac{z}{\vert{}\mathbf{r}\vert{}} \end{bmatrix}$$

where $R_E$ is the Earth's equatorial radius, $\mu$ is the standard gravitational parameter, and $x, y, z$ are the ECI coordinates of the vehicle. This addition is computationally cheap but elevates the trajectory propagation to true orbital fidelity.

\subsection{Exo-Atmospheric Thermal Equilibrium}
While the structural vibrations damp out, the thermal solver remains active but undergoes a boundary condition inversion. With convective heating gone ($\mathbf{Q}_{thermal} = \mathbf{0}$), the exterior Neumann boundary condition of the 2D spectral plate is governed purely by Stefan-Boltzmann radiative cooling:

$$-k_{thermal} \frac{\partial T}{\partial n} = -\epsilon \sigma T^4$$

During the long fractional orbit, the vehicle acts as a blackbody radiator. The massive heat load absorbed during the boost phase is slowly radiated away into deep space, while internal conduction diffuses the thermal gradients through the thickness map $m(\xi, \eta)$. By the time the FOBS commands its de-orbit burn, the thermal solver will have accurately tracked the vehicle's cool-down, resetting the initial state for the extreme hypersonic reentry phase.

\subsection{The De-Orbit Burn Kinematics}
To terminate the fractional orbit, the vehicle must execute a retrograde burn. This re-engages the 6-DOF quaternion kinematics.

The flight computer commands a target attitude $\mathbf{q}_{target}$ to orient the vehicle's main engine against the velocity vector. A standard Proportional-Derivative (PD) quaternion controller commands the body-frame angular velocity $\boldsymbol{\omega}_B$ to align the vehicle:

$$\mathbf{T}_{burn} = \text{diag}(T_x, T_y, T_z)$$

where $T_x$ is the primary retro-thrust.

The translational kinematics equation briefly re-awakens the RHS forcing vector:

$$\mathbf{\dot{v}}_E = \mathbf{g}(\mathbf{r}) + \frac{1}{m_{bulk}} \mathbf{T}_I^B(\mathbf{q})^T \mathbf{T}_{burn}$$

Once the required $\Delta V$ is achieved to drop the perigee into the sensible atmosphere, the engine cuts off. As the vehicle descends, $\rho(h)$ begins to rise, the dynamic pressure activates, and the RKF45 solver automatically tightens its time step as the Newtonian pressure and Sutton-Graves heating smoothly return to dominate the physics engine.

This architecture allows the framework to transition from boost, to orbital coast, to hypersonic reentry seamlessly within a single, unbroken numerical integration run.

\section{Conclusion}
\label{sec:conclusion}
The development of the PASSES-HGV framework effectively bridges the critical gap between rapid conceptual trajectory modeling and the high-fidelity, multi-physics simulations required for production-level defense architectures. By systematically dismantling the limiting assumptions of legacy engineering solvers, this framework provides a comprehensive, tightly coupled ecosystem for hypersonic glide vehicle analysis.

The structural and thermodynamic domains have been elevated to capture the true 3D aeroelastic and extreme heating realities of lifting bodies operating in the hypersonic regime. Replacing thin-plate approximations with a continuous Mindlin-Reissner anisotropic formulation allows for rigorous transverse shear analysis along varying waverider cross-sections. Simultaneously, the integration of Fay-Riddell and Tauber-Sutton correlations ensures that finite-rate chemistry, surface catalysis, and volumetric radiative heating are evaluated dynamically, providing physically accurate boundary conditions for the thermal protection system without breaking the continuous integration loop.

In the guidance and navigation domains, transitioning to Successive Convexification transforms autonomous trajectory generation into a strictly bounded optimal control problem, mathematically guaranteeing thermal and dynamic constraint satisfaction. Coupling this with an ionization-aware Extended Kalman Filter ensures the flight computer can autonomously manage RF blackout windows. By dynamically sensing the plasma sheath and transitioning from GNSS reliance to pure inertial propagation, the framework seamlessly manages state estimation within the unforgiving margins of the flight corridor.

Ultimately, PASSES-HGV represents a fundamental advancement in applied mathematics and advanced modeling and simulation. By casting the entire multi-physics environment as a series of continuous block-Kronecker tensor operations, the framework leverages GPU acceleration to execute massive batched Monte Carlo statistical validations. This architecture delivers the precise probabilistic rigor, deterministic control, and computational efficiency demanded by modern applied physics laboratories and national research centers, providing a definitive, high-performance toolset for the next generation of hypersonic defense engineering.

% \section*{Acknowledgments}

\pagebreak

\printbibliography

\pagebreak

\begin{appendices}

\end{appendices}

\end{document}
